% ══════════════════════════════════════════════════════════════════════
% Appendix A  Mixed-rank SU-SU no-Veneziano classes
% ══════════════════════════════════════════════════════════════════════

In the main text we focused on the equal-rank case
$\SU(N)_{0}\times\SU(N)_{1}$ with rank multipliers $(m_{0},m_{1})=(1,1)$.
The same classification pipeline may be run at general coprime rank
multipliers, producing additional non-Veneziano universality classes
for
\begin{equation}
  G\;=\;\SU(m_{0}N)_{0}\times\SU(m_{1}N)_{1},
  \qquad\gcd(m_{0},m_{1})=1\,,
\end{equation}
with $(m_{0},m_{1})$ chosen up to a swap of the two nodes.  We retain
only classes with $m_{0},m_{1}\le 4$.

\paragraph{Inventory.}  Table~\ref{tab:mixed-rank-inventory} lists the
$20$ mixed-rank non-Veneziano classes with $a_{/N^{2}}$ defined.
As in the equal-rank case, every non-degenerate class saturates
$a=c$ to leading order in $1/N$.

\begin{table}[h]
  \centering
  \small
  \begin{tabular}{cccrll}
    \toprule
    Class & $m_{0}$ & $m_{1}$ & \# theories & $a/N^{2}$ & $a/c$ \\
    \midrule
    $111$ & $2$ & $1$ & $2$  & $-0.6328$ & $1.0$ \\
    $117$ & $2$ & $1$ & $2$  & $0$ & (deg.) \\
    $121$ & $2$ & $1$ & $8$  & $0.7910$ & $1.0$ \\
    $122$ & $2$ & $1$ & $8$  & $0.7922$ & $1.0$ \\
    $134$ & $2$ & $1$ & $8$  & $1.0547$ & $1.0$ \\
    $136$ & $2$ & $1$ & $16$ & $1.0760$ & $1.0$ \\
    $153$ & $2$ & $1$ & $20$ & $1.1470$ & $1.0$ \\
    $159$ & $2$ & $1$ & $3$  & $1.1897$ & $1.0$ \\
    $160$ & $2$ & $1$ & $40$ & $1.1926$ & $1.0$ \\
    $174$ & $2$ & $1$ & $6$  & $1.25$ & $1.0$ \\
    $283$ & $3$ & $1$ & $8$  & $0.9722$ & $1.0$ \\
    $298$ & $3$ & $1$ & $20$ & $2.3056$ & $1.0$ \\
    $377$ & $3$ & $2$ & $8$  & $2.2222$ & $1.0$ \\
    $384$ & $3$ & $2$ & $34$ & $2.4668$ & $1.0$ \\
    $398$ & $3$ & $2$ & $20$ & $2.5477$ & $1.0$ \\
    $437$ & $3$ & $2$ & $12$ & $2.9722$ & $1.0$ \\
    $450$ & $3$ & $2$ & $88$ & $3.1192$ & $1.0$ \\
    $681$ & $4$ & $3$ & $34$ & $4.9951$ & $1.0$ \\
    $730$ & $4$ & $3$ & $12$ & $5.9425$ & $1.0$ \\
    $731$ & $4$ & $3$ & $88$ & $5.9568$ & $1.0$ \\
    \bottomrule
  \end{tabular}
  \caption{Mixed-rank non-Veneziano $\SU(m_{0}N)\times\SU(m_{1}N)$
    universality classes with $m_{0}\ne m_{1}$ (or $m_{0}=m_{1}>1$).
    $a/N^{2}$ is quoted to four decimal places; exact forms are
    available in the accompanying database.}
  \label{tab:mixed-rank-inventory}
\end{table}

\paragraph{RG phase distribution.}  Running the boundary analysis of
Section~\ref{sec:boundary} on each mixed-rank class reveals the same
three topology types (i), (ii), (iii) as in the equal-rank sector.
As in Section~\ref{sec:boundary-phases}, \emph{none} of the
non-Veneziano mixed-rank classes realise topology (iii) at $\Nf=0$.
Single-parameter deformations analogous to the class-$16$ deformation
of Section~\ref{sec:phases-class16} can push individual classes into
the Figure~6 window; characterising which mixed-rank classes admit a
Figure~6 phase is a straightforward but tedious exercise we leave for
future work.

\paragraph{Negative $a/N^{2}$.}  Class $111$ has a negative value of
the $a$-maximum at leading order, which violates the expected
positivity of $a$ for a unitary SCFT.  The negative value signals
that the $a$-max critical point at this order is not within the
physical R-symmetry polytope, and the theory is either non-unitary or
its leading-$N$ analysis misses a relevant sub-leading correction.
We flag this class as a cautionary boundary case rather than including
it in the topology-classification tables.
